\documentclass[10pt]{article}
\usepackage[margin=1in]{geometry}
\usepackage{amsmath,amssymb,amsthm,hyperref}
\newtheorem{theorem}{Theorem}
\newtheorem{lemma}{Lemma}
\newtheorem{corollary}{Corollary}
\DeclareMathOperator{\lcm}{lcm}
\title{MathDB 367325: polynomial blocks for every\ Motzkin convolution power}
\author{Principia Math In Progress}
\date{18 September 2026}
\begin{document}
\maketitle

\section*{Task 1: Full proof of the determinant block pattern}
Let $M(z)=\sum_{n\ge0}M_nz^n$ be the Motzkin generating function,
so $M=1+zM+z^2M^2$. For an integer $r\ge1$, write
\[
 M(z)^r=\sum_{n\ge0}a_n^{(r)}z^n,\qquad
 H_N^{(r)}=\det(a_{i+j}^{(r)})_{0\le i,j<N},\qquad H_0^{(r)}=1.
\]
Wang and Zhang's Conjecture 8
(\href{https://arxiv.org/abs/2502.21050}{arXiv:2502.21050}, p.~4)
predicts shifted block lengths $r$ when $3\mid r$ and $3r$ otherwise.
We establish the determinant-level assertion in the precise form
below: each residue class is a polynomial in the block index, with
an explicitly determined parity factor. These are the polynomial
patterns displayed in their Theorems 1, 2, 4--6. The block length
here is not the number of continued-fraction transformations: for
example the source uses two transformations for block length three
and ten for block length twelve. No minimality of a block length
or periodicity of the numerical values themselves is required.

\begin{theorem}\label{main}
Set
\[
 p=\lcm(r,3)=\begin{cases}r,&3\mid r,\\3r,&3\nmid r,\end{cases}
 \qquad \epsilon=(-1)^{\binom r2},\qquad
 \kappa=\begin{cases}r,&3\mid r,\\r-1,&3\nmid r,\end{cases}
 \qquad D=\binom\kappa2.
\]
For each $s\in\{0,\ldots,p-1\}$ there is an integer-valued polynomial
$B_{r,s}(X)\in\mathbb Q[X]$, of degree at most $D$, such that
\begin{equation}\label{blocks}
 H_{pn+s}^{(r)}=\epsilon^n B_{r,s}(n)\qquad(n\ge0).
\end{equation}
The proof supplies an exact determinant of order $2r-2$ for every
$H_N^{(r)}$, as well as a finite formula for every block polynomial.
The convention is $D=0$ when $r=1$.
\end{theorem}

Define monic shifted Chebyshev polynomials by
\[
 P_{-1}(t)=0,\quad P_0(t)=1,\quad P_1(t)=t-1,\qquad
 P_{j+1}(t)=(t-1)P_j(t)-P_{j-1}(t).
\]
Thus $P_j(1+2\cos\theta)=\sin((j+1)\theta)/\sin\theta$.
Let $\mu$ be the probability measure
\[
 d\mu(t)=\frac{\sqrt{4-(t-1)^2}}{2\pi}\,dt\quad(-1\le t\le3).
\]
The substitution $t=1+2\cos\theta$ gives
\begin{equation}\label{orthogonal}
 \int P_i(t)P_j(t)\,d\mu(t)
 =\frac2\pi\int_0^\pi\sin((i+1)\theta)\sin((j+1)\theta)\,d\theta
 =\delta_{ij}.
\end{equation}
Its moments are $M_n$: expansion of $(1+2\cos\theta)^n$ gives
\[
 \int t^n\,d\mu(t)
 =\sum_{k=0}^{\lfloor n/2\rfloor}\binom n{2k}\frac1{k+1}\binom{2k}k=M_n.
\]
Here the odd centered moments vanish by symmetry and the even ones
are the displayed Catalan numbers, by the elementary sine integral.

\begin{lemma}[The weight for a convolution power]\label{weight}
For $Q_r(t)=t^{r-1}P_{r-1}(t)$ and every $n\ge0$,
\begin{equation}\label{moments}
 a_n^{(r)}=\int t^n Q_r(t)\,d\mu(t).
\end{equation}
\end{lemma}
\begin{proof}
In formal Laurent series at infinity, put
$G(w)=w^{-1}M(w^{-1})=\sum_{n\ge0}M_nw^{-n-1}$.
The Motzkin equation gives $G^2=(w-1)G-1$. Induction using the
recurrence for $P_j$ consequently gives
\[
 G(w)^r=P_{r-1}(w)G(w)-P_{r-2}(w).
\]
Multiply by $w^{r-1}$. The left side is $w^{-1}M(w^{-1})^r$.
The second term on the right is a polynomial, and the coefficient
of $w^{-n-1}$ in the first is
$\sum_j[t^j]Q_r(t)\,M_{n+j}$, proving \eqref{moments}.
\end{proof}

We recall and prove the needed form of the classical
Christoffel determinant identity. It is also Theorem 1 of
Krattenthaler, \href{https://arxiv.org/abs/2101.04225}{arXiv:2101.04225},
with confluent evaluation as in its Proposition 5; the identity itself
is an established tool.

\begin{lemma}[Polynomial modification of moments]\label{christoffel}
Let $Q(t)=\prod_{i=1}^d(t-\alpha_i)$ be monic and let the monic
$P_j$ be orthonormal for $\mu$. For distinct roots,
\begin{equation}\label{distinct}
 \det\left(\int t^{i+j}Q(t)\,d\mu(t)\right)_{0\le i,j<N}
 =(-1)^{Nd}\frac{\det(P_{N+j}(\alpha_i))_{1\le i\le d,\,0\le j<d}}
 {\det(\alpha_i^j)_{1\le i\le d,\,0\le j<d}}.
\end{equation}
Repeated roots are handled by replacing the rows at a root of
multiplicity $m$ by derivative evaluations of orders $0,\ldots,m-1$,
each divided by the factorial of its order, in both determinants.
The confluent denominator is nonzero.
\end{lemma}
\begin{proof}
Expansion of two determinants and integration term by term give
the left side as
\[
 \frac1{N!}\int\cdots\int
 \Delta(x_1,\ldots,x_N)^2\prod_{k=1}^N Q(x_k)\prod_{k=1}^N d\mu(x_k),
\]
where $\Delta$ uses increasing powers in its columns. The mixed
Vandermonde, with the $x$ variables followed by the $\alpha$ variables,
satisfies
\[
 \Delta(x,\alpha)=\Delta(x)\Delta(\alpha)\prod_{k,i}(\alpha_i-x_k).
\]
Replace its monomials, and those in the remaining $\Delta(x)$,
by the monic $P_j$, which does not change either determinant.
Expand the mixed determinant along its last $d$ rows. By
orthonormality, integration against the other determinant kills
every term except that using columns $0,\ldots,N-1$ for the first
$N$ rows. That term has positive expansion sign and integral $N!$.
The remaining minor is the numerator in \eqref{distinct}, and
$\prod Q(x_k)$ contributes $(-1)^{Nd}$.

For colliding roots, divide numerator and denominator by their
common within-cluster Vandermonde factors and take limits. Taylor
expansion gives exactly the stated derivative rows. These derivative
evaluations are linearly independent on polynomials of degree less
than $d$: a polynomial in their common kernel is divisible by $Q$
and hence is zero. The confluent denominator is therefore nonzero.
Both sides depend continuously on the polynomial coefficients, so
the identity remains valid. The empty-determinant case is included.
\end{proof}

\begin{proof}[Proof of Theorem~\ref{main}]
For $r=1$, \eqref{orthogonal} gives $H_N^{(1)}=1$, which proves
the assertion. Suppose $r\ge2$ and put $d=2r-2$.
The roots of $P_{r-1}$ are simple and are
\[
 \rho_j=1+2\cos(j\pi/r),\qquad 1\le j\le r-1.
\]
One equals zero exactly when $3\mid r$, in which case its index is
$j_0=2r/3$, an even integer. Thus $Q_r$ has a root zero of
multiplicity $\kappa$, together with the nonzero $\rho_j$, all simple.

Let $\mathcal E_N$ be the $d\times d$ matrix whose columns are
$P_N,P_{N+1},\ldots,P_{N+d-1}$ and whose rows apply the following
functionals: $f\mapsto f^{(k)}(0)/k!$ for $0\le k<\kappa$, then
$f\mapsto f(\rho_j)$ for each nonzero $\rho_j$, in any fixed order.
Let $\mathcal V$ use the same rows and columns $1,t,\ldots,t^{d-1}$.
Lemmas~\ref{weight} and \ref{christoffel}, and evenness of $d$, give
\begin{equation}\label{small}
 H_N^{(r)}=\frac{\det\mathcal E_N}{\det\mathcal V}.
\end{equation}
In particular no division by a possibly vanishing $H_N^{(r)}$ occurs.

Set $N=pn+s$. The integer $p/r$ is odd. For a nonzero $\rho_j$,
the sine formula gives, for every column offset $h$,
\[
 P_{pn+s+h}(\rho_j)=(-1)^{nj}P_{s+h}(\rho_j).
\]
Factor this sign out of each such row. Their product is
$(-1)^{n\binom r2}=\epsilon^n$: the only possibly omitted index
is the even integer $j_0$.

For the zero rows use $t=1+2\cos\theta$ near $\theta_0=2\pi/3$.
Differentiation in $t$ acts as
$-(2\sin\theta)^{-1}d/d\theta$. After $k$ differentiations of
$\sin((L+1)\theta)/\sin\theta$, the value at $\theta_0$ is a
linear combination of
\[
 (L+1)^a\sin((L+1)\theta_0),\qquad
 (L+1)^a\cos((L+1)\theta_0),\qquad 0\le a\le k,
\]
with coefficients independent of $L$. This follows directly by
induction with the product rule; $\sin\theta_0\ne0$ ensures all
coefficients are finite. Since $3\mid p$, substituting $L=pn+s+h$
makes the sine and cosine factors independent of $n$.
Each entry in the derivative row of order $k$ is therefore a
polynomial in $n$ of degree at most $k$.

After removing the simple-root signs, the determinant in
\eqref{small} is a polynomial in $n$ of degree at most
$0+1+\cdots+(\kappa-1)=D$. Its denominator is a fixed nonzero
constant, so this proves \eqref{blocks} initially over $\mathbb R$.
Its values at all nonnegative integers are integers, because both
$H_N^{(r)}$ and $\epsilon$ are integers. Newton interpolation now
gives the explicit expression
\begin{equation}\label{newton}
 B_{r,s}(X)=\sum_{j=0}^{D}\binom Xj
 \sum_{h=0}^{j}(-1)^{j-h}\binom jh\epsilon^h H_{ph+s}^{(r)}.
\end{equation}
Hence $B_{r,s}\in\mathbb Q[X]$ and is integer-valued (also at
negative integers).

For completeness, \eqref{small} has a version requiring only
integer polynomial arithmetic. Reduce $P_{N+j}$ modulo the monic
polynomial $Q_r$ and write its remainder as $\sum_{i=0}^{d-1}c_{ij}t^i$.
Every evaluation functional defining $\mathcal E_N$ annihilates
multiples of $Q_r$, so $\mathcal E_N=\mathcal V(c_{ij})$ and
\begin{equation}\label{integer}
 H_N^{(r)}=\det(c_{ij})_{0\le i,j<2r-2}.
\end{equation}
This proves the promised finite determinant formula and completes
the proof for every exponent, residue, and nonnegative block index.
\end{proof}

\begin{corollary}
The determinant sequence satisfies
$(E^p-\epsilon)^{D+1}H^{(r)}=0$, where $(EH)_N=H_{N+1}$.
In particular $\sum_{N\ge0}H_N^{(r)}z^N$ is rational, with
denominator dividing $(1-\epsilon z^p)^{D+1}$.
\end{corollary}
\begin{proof}
On each residue class, applying $E^p-\epsilon$ takes a polynomial
to $\epsilon$ times its forward difference. After $D+1$ applications
it vanishes. Multiplication of the ordinary generating series by
the displayed denominator therefore leaves a polynomial.
\end{proof}

\paragraph{Attribution and scope.}
The source authors conjectured the block lengths and provided
explicit cases. Polynomial modifications of orthogonal-polynomial
moments and their Christoffel determinant formula are classical;
the latter is proved here for completeness and attributed above.
The contribution here is the convolution weight $Q_r$, the analysis
of its roots, and the resulting uniform signed polynomial blocks.
This addresses Conjecture 8 for the determinants; it does not assert
the separate numerical formulas of Conjectures 3 and 7 or a specific
orbit of the continued-fraction transformation. The displayed sign
is derived independently, rather than imported from those formulas.
Exactly one MathDB problem is addressed. The following tasks are
proposed directions, not additional solved problems.

\section*{Task 2: Find sharp degrees and leading coefficients}
Determine the exact degree of every nonzero $B_{r,s}$, classify the
identically zero residues, and evaluate the leading coefficients.
The general bound $\binom\kappa2$ need not be sharp.

\section*{Task 3: Evaluate distinguished residue classes}
Use the fixed-size determinant to prove explicit formulas for
$B_{r,0}$, $B_{r,1}$, and the paired residues in the source's
Conjectures 3 and 7, including their correct sign factors.

\section*{Task 4: Determine the least block length}
Classify all cancellations that allow a smaller signed polynomial
block length or a smaller generating-function denominator than
those supplied by the theorem.

\section*{Task 5: Extend to weighted Motzkin moments}
Replace the horizontal-step weight one by a real parameter.
Determine which parameters produce polynomial blocks, and describe
the exponential-polynomial behavior when the relevant root angles
are not rational multiples of $\pi$.
\end{document}
